\scanpage{433}
\mspage{345}

l'implication $a\Longrightarrow c$ (via $a\Rightarrow a'\overset{!}{\Rightarrow}c$). Donc Il faut surtout voir ce théorème comme donnant un critère remarquable d'appartenance à $\mathrm{Ob}\,\underline S$ d'un objet $x$ de $\unsure{M|\mathcal M}$ qui est présenté par un syst. inductif $\unsure{\Phi|\mathscr I|\mathscr D}$-$\pi$-adapté :
\[
x\unsure{+|\in}\mathrm{Ob}\,\underline S\quad\Longleftrightarrow\quad I_{\underline S}\text{ cofinal},\tag*{\unsure{$(*)$|$(**)$}}
\]
et dans cette équivalence, c'est l'implication $\Longrightarrow$ qui est la partie \deleted{``vache''} \inserted{non-triviale}, et remarquable.

\uline{Proposition} 4.13.\deleted{\unsure{23}}24. Soient \deleted{\unsure{$\Phi$}}$M$, $\pi$, $\Phi$ comme dans 4.13.23, \unsure{$*$} on suppose que \deleted{$M$ \unsure{stable}} la condition \uline{1°} de \deleted{\unsure{la dé}} 4.13.23 est satisfaite ($\pi$ utile pour $M$, $M$ stable par $\varinjlim$ de type $\Phi$).

% slanted
\margin{NB. La condition $I_{\alpha_0}\in\Phi$ n'est pas utilisée ici.}

Soit $\Sigma_\pi^{\prime\Phi}\deleted{=\unsure{\Sigma(M,\pi,\Phi)}}$ l'ens des sous-catégories \deleted{que} \inserted{$\underline S$} de $M$ qui sont \deleted{\unsure{$\underline S$}} $\Phi$-$\pi$-accessibles, dans $M$, et $\Sigma^{\Phi}_\pi$ l'ens des sous-catégories \inserted{$\underline S(\pi)$} de $M_\pi$ qui sont str. pleines, \deleted{\unsure{Karou}} et stables dans $M_\pi$ par $\varinjlim$ de type $\Phi$. (Ces sous-catégories sont forcément Karoubiennes, car \deleted{$I_{\alpha_0}$} stables par $\varinjlim$ de type $I_{\alpha_0}$, donc aussi par facteurs directs). Alors les bijections (4.13.\deleted{\unsure{1}}4.1) (p. \deleted{\unsure{p}} 288) induisent des bijections inverses l'une de l'autre
\diagram{\begin{tikzcd}[ampersand replacement=\&]
 \Sigma_{\pi}^{\Phi}\arrow[r,shift left] \& {\Sigma_{\pi}^{\prime\Phi}.}\arrow[l,shift left]
\end{tikzcd}}

% archive sheet number: 432

